\chapter{Alan Baker --- Fields Medal 1970}\label{ch1970alanbaker:alan-baker-fields-medal-1970-a-comprehensive-report}

\begin{center}\rule{0.5\linewidth}{0.5pt}\end{center}

\textbf{Prepared:} July 2026 \textbf{Subject:} Alan Baker \textbf{Born:}
August 19, 1939, London, England \textbf{Died:} February 4, 2018
\textbf{Position:} Professor, University of Cambridge \textbf{Doctoral
Advisor:} Harold Davenport (University of Cambridge)

\textbf{Fields Medal Citation:} \emph{``For his work on transcendental
number theory, in particular for his method of linear forms in
logarithms of algebraic numbers, which has revolutionized that field and
led to decisive advances in several classical problems of Diophantine
geometry.''}

\begin{center}\rule{0.5\linewidth}{0.5pt}\end{center}

\section{1. Introduction}\label{ch1970alanbaker:introduction}

Alan Baker was awarded the Fields Medal in 1970 for his substantial
contributions to number theory.

His \textbf{theory of linear forms in logarithms} provided a powerful
new tool for solving Diophantine equations, leading to effective bounds
for solutions to many classical problems.

Baker's method transformed transcendental number theory from a subject
of existence proofs into a domain of explicit effective computation.

His results include: - Effective bounds for linear forms in logarithms.
- The effective solution of Thue's equation. - The effective solution of
the Mordell equation. - The effective solution of hyperelliptic
equations. - Results on Catalan's equation. - Results on the class
number of imaginary quadratic fields. - Transcendence of values of
\(L\)-functions.

\begin{center}\rule{0.5\linewidth}{0.5pt}\end{center}

\section{2. Biographical Background}\label{ch1970alanbaker:biographical-background}

\subsection{2.1 Early Life and
Education}\label{ch1970alanbaker:early-life-and-education}

Alan Baker was born on August 19, 1939, in London, England.

He attended the University of Cambridge (Trinity College) as an
undergraduate.

He earned his Ph.D.~in 1964 under Harold Davenport, with a thesis on
Diophantine approximation.

\subsection{2.2 Academic Career}\label{ch1970alanbaker:academic-career}

\begin{itemize}
\tightlist
\item
 \textbf{1958--1964:} Undergraduate and Ph.D.~student, University of
 Cambridge (Trinity College).
\item
 \textbf{1964--1967:} Research Fellow, Trinity College, Cambridge.
\item
 \textbf{1967--1968:} Visiting Professor, University of Michigan, Ann
 Arbor.
\item
 \textbf{1968--1974:} Lecturer, then Reader, University of Cambridge.
\item
 \textbf{1970:} Fields Medal at the ICM in Nice.
\item
 \textbf{1974--2004:} Professor of Pure Mathematics, University of
 Cambridge.
\item
 \textbf{2004:} Professor Emeritus, Cambridge.
\end{itemize}

\subsection{2.3 Doctoral Work}\label{ch1970alanbaker:doctoral-work}

Baker's early work on Diophantine approximation led him to develop
effective methods for bounding solutions to Diophantine equations.

His Ph.D.~thesis already contained the seeds of his method of linear
forms in logarithms.

\subsection{2.4 Influence of Davenport}\label{ch1970alanbaker:influence-of-davenport}

Harold Davenport was a leading figure in analytic number theory.

He guided Baker toward the problems of Diophantine approximation and the
theory of exponential sums.

Davenport's emphasis on explicit bounds shaped Baker's approach to
number theory.

\begin{center}\rule{0.5\linewidth}{0.5pt}\end{center}

\section{3. Linear Forms in
Logarithms}\label{ch1970alanbaker:linear-forms-in-logarithms}

\subsection{3.1 Statement of the
Theorem}\label{ch1970alanbaker:statement-of-the-theorem}

Let \(\alpha_1, \dots, \alpha_n\) be algebraic numbers and
\(\beta_1, \dots, \beta_n\) be rational integers.

Consider the linear form

\[
\Lambda = \beta_1 \log \alpha_1 + \cdots + \beta_n \log \alpha_n.
\]

Baker's theorem provides an effective lower bound: if
\(\Lambda \neq 0\), then

\[
|\Lambda| > H^{-C}
\]

where \(H = \max(|\beta_1|, \dots, |\beta_n|)\) and \(C\) is a constant
depending on \(\alpha_1, \dots, \alpha_n\).

\subsection{3.2 The Algebraic Version}\label{ch1970alanbaker:the-algebraic-version}

For algebraic numbers \(\alpha_i\) and integers \(\beta_i\), the
quantity

\[
\Lambda = \alpha_1^{\beta_1} \cdots \alpha_n^{\beta_n} - 1
\]

also satisfies a nontrivial lower bound when nonzero.

This is the form most commonly used in applications to Diophantine
equations.

\subsection{3.3 The Method of Proof}\label{ch1970alanbaker:the-method-of-proof}

Baker's proof combines several techniques: 1. \textbf{Gelfond's method}
for linear forms in two logarithms. 2. \textbf{Padé approximation} and
the construction of auxiliary functions. 3. \textbf{Hensel's lemma} and
\(p\)-adic methods. 4. \textbf{The theory of algebraic number fields}
and heights of algebraic numbers. 5. \textbf{Complex analysis} and the
maximum modulus principle.

The auxiliary function is constructed as a linear combination of
exponentials with carefully chosen algebraic coefficients.

A delicate interpolation argument at many points yields a contradiction
if the linear form is too small.

\subsection{3.4 Refinements and
Generalizations}\label{ch1970alanbaker:refinements-and-generalizations}

Baker's initial results have been significantly improved: - Masser's
vanishing theorem. - The Birch--Masser theorem. - Waldschmidt's sharper
bounds. - Mihailescu's use of algebraic number theory in Catalan's
conjecture.

\begin{center}\rule{0.5\linewidth}{0.5pt}\end{center}

\section{4. Effective Solution of Thue's
Equation}\label{ch1970alanbaker:effective-solution-of-thues-equation}

\subsection{4.1 Statement}\label{ch1970alanbaker:statement}

A \textbf{Thue equation} is a Diophantine equation

\[
F(x,y) = k
\]

where \(F\) is a homogeneous form of degree \(d \geq 3\) with integer
coefficients and \(k \in \mathbb{Z}\).

Thue (1909) proved that such equations have only finitely many integer
solutions, but his proof was non-effective.

\subsection{4.2 Baker's Theorem}\label{ch1970alanbaker:bakers-theorem}

Baker proved: for any Thue equation, the integer solutions satisfy

\[
\max(|x|, |y|) < \exp(C)
\]

where \(C\) is an effective computable constant depending on \(F\) and
\(k\).

This makes it possible in principle to find all solutions by finite
search.

\subsection{4.3 Impact}\label{ch1970alanbaker:impact}

Baker's result effectively solves the Thue equation, a problem open
since 1909.

The method has been implemented computationally, allowing explicit
determination of all solutions.

\begin{center}\rule{0.5\linewidth}{0.5pt}\end{center}

\section{5. Effective Solution of the Mordell
Equation}\label{ch1970alanbaker:effective-solution-of-the-mordell-equation}

\subsection{5.1 Statement}\label{ch1970alanbaker:statement-1}

The \textbf{Mordell equation} is

\[
y^2 = x^3 + k
\]

for integer \(k \neq 0\).

Mordell (1922) proved that the rational points on elliptic curves form a
finitely generated group, but this did not provide an effective bound
for integer solutions.

\subsection{5.2 Baker's Theorem}\label{ch1970alanbaker:bakers-theorem-1}

Baker proved: the integer solutions of \(y^2 = x^3 + k\) satisfy

\[
\max(|x|, |y|) < \exp(\exp(C))
\]

where \(C\) is effective and depends only on \(k\).

This reduces the search for all integer points to a finite computation.

\subsection{5.3 The Class Number
Problem}\label{ch1970alanbaker:the-class-number-problem}

Baker's method also applies to the determination of imaginary quadratic
fields with class number 1.

He proved that if \(\mathbb{Q}(\sqrt{-d})\) has class number 1, then
\(d\) belongs to the finite set

\[
\{1, 2, 3, 7, 11, 19, 43, 67, 163\}.
\]

This was proved simultaneously by Stark (1967) and completed the work of
Heegner (1952).

\begin{center}\rule{0.5\linewidth}{0.5pt}\end{center}

\section{6. The Hyperelliptic
Equation}\label{ch1970alanbaker:the-hyperelliptic-equation}

\subsection{6.1 Statement}\label{ch1970alanbaker:statement-2}

A \textbf{hyperelliptic equation} (genus \(\geq 1\)) is

\[
y^2 = f(x)
\]

where \(f(x)\) is a polynomial with integer coefficients and at least
three distinct roots.

\subsection{6.2 Baker's Result}\label{ch1970alanbaker:bakers-result}

Baker proved effective bounds for integer solutions of hyperelliptic
equations.

The method uses the factorization

\[
y - \sqrt{f(x)} = P(x) + Q(x)\sqrt{f(x)}
\]

and bounds the resulting linear forms in logarithms.

\subsection{6.3 Applications}\label{ch1970alanbaker:applications}

\begin{itemize}
\tightlist
\item
 The equation \(y^2 = x^n + k\) for fixed \(n \geq 3\).
\item
 The equation \(y^3 = x^2 + k\).
\item
 The generalized Ramanujan--Nagell equation.
\end{itemize}

\begin{center}\rule{0.5\linewidth}{0.5pt}\end{center}

\section{7. The Catalan Equation}\label{ch1970alanbaker:the-catalan-equation}

\subsection{7.1 Statement}\label{ch1970alanbaker:statement-3}

\textbf{Catalan's equation} (conjectured 1844) is

\[
x^p - y^q = 1
\]

where \(x, y > 0\) and \(p, q > 1\).

Catalan conjectured that the only solution is

\[
3^2 - 2^3 = 1\quad (x=3, y=2, p=2, q=3).
\]

\subsection{7.2 Baker's Contribution}\label{ch1970alanbaker:bakers-contribution}

Baker proved effective bounds for \(x, y, p, q\) using linear forms in
logarithms.

Specifically: - The solutions satisfy \(\max(p,q) \leq C\) for some
absolute constant \(C\). - The solutions satisfy

\[
\max(x,y) \leq \exp(\exp(C)).
\]

This reduced Catalan's conjecture to a finite computation.

\subsection{7.3 Mihailescu's Proof}\label{ch1970alanbaker:mihailescus-proof}

Mihailescu (2002) used Baker's method together with deeper algebraic
number theory to prove the full conjecture.

This was one of the most well-known applications of Baker's theory.

\begin{center}\rule{0.5\linewidth}{0.5pt}\end{center}

\section{8. Transcendence Theory}\label{ch1970alanbaker:transcendence-theory}

\subsection{8.1 The Gelfond--Schneider
Theorem}\label{ch1970alanbaker:the-gelfondschneider-theorem}

The \textbf{Gelfond--Schneider theorem} (1934) solved Hilbert's 7th
problem:

If \(\alpha\) is algebraic, \(\alpha \neq 0, 1\), and \(\beta\) is
algebraic irrational, then \(\alpha^\beta\) is transcendental.

\subsection{8.2 Baker's Generalization}\label{ch1970alanbaker:bakers-generalization}

Baker proved that linear combinations of logarithms of algebraic numbers
with algebraic coefficients are transcendental, unless the combination
vanishes.

If \(\beta_0, \beta_1, \dots, \beta_n\) are algebraic with
\(\beta_0 \neq 0\) or \(\beta_1, \dots, \beta_n\) not all zero, then

\[
\beta_0 + \beta_1 \log \alpha_1 + \cdots
+ \beta_n \log \alpha_n \neq 0.
\]

\subsection{8.3 Algebraic Independence}\label{ch1970alanbaker:algebraic-independence}

Baker's methods led to results on \textbf{algebraic independence}:
values of the exponential function at algebraically independent points
are themselves algebraically independent.

\begin{center}\rule{0.5\linewidth}{0.5pt}\end{center}

\section{9. The Mordell--Weil Theorem}\label{ch1970alanbaker:the-mordellweil-theorem}

\subsection{9.1 Background}\label{ch1970alanbaker:background}

The \textbf{Mordell--Weil theorem} states that the set of rational
points on an elliptic curve

\[
E(\mathbb{Q})
\]

is a finitely generated abelian group.

The rank is the number of independent generators.

\subsection{9.2 Baker's Effective
Version}\label{ch1970alanbaker:bakers-effective-version}

Baker provided effective bounds for the height of generators of
\(E(\mathbb{Q})\), using linear forms in elliptic logarithms.

This makes the computation of Mordell--Weil groups feasible in many
cases.

\begin{center}\rule{0.5\linewidth}{0.5pt}\end{center}

\section{\texorpdfstring{10. \(p\)-adic
Methods}{10. p-adic Methods}}\label{ch1970alanbaker:p-adic-methods}

\subsection{10.1 Background}\label{ch1970alanbaker:background-1}

Baker's method was extended to the \(p\)-adic setting using \(p\)-adic
logarithms and exponentials.

The \(p\)-adic linear forms in logarithms provide effective bounds for
solutions to \(p\)-adic Diophantine equations.

\subsection{10.2 Applications}\label{ch1970alanbaker:applications-1}

\begin{itemize}
\tightlist
\item
 The \(p\)-adic solution of Thue equations.
\item
 Bounds for \(S\)-unit equations.
\item
 The \(p\)-adic Mordell equation.
\end{itemize}

\begin{center}\rule{0.5\linewidth}{0.5pt}\end{center}

\section{\texorpdfstring{11. The \(abc\)
Conjecture}{11. The abc Conjecture}}\label{ch1970alanbaker:the-abc-conjecture}

\subsection{11.1 Statement}\label{ch1970alanbaker:statement-4}

The \textbf{\(abc\) conjecture} (Masser--Oesterlé) states: for
\(\epsilon > 0\), there exists \(C(\epsilon)\) such that for coprime
integers \(a+b=c\),

\[
\max(|a|, |b|, |c|) \leq C(\epsilon)
\cdot \text{rad}(abc)^{1+\epsilon}
\]

where \(\text{rad}(n)\) is the product of distinct primes dividing
\(n\).

\subsection{11.2 Baker's Work}\label{ch1970alanbaker:bakers-work}

Baker proved results relating linear forms in logarithms to the \(abc\)
conjecture.

He showed that effective versions of the \(abc\) conjecture would follow
from sharpened bounds for linear forms in logarithms.

\begin{center}\rule{0.5\linewidth}{0.5pt}\end{center}

\section{12. Awards and Recognition}\label{ch1970alanbaker:awards-and-recognition}

\textbf{1970:} - Fields Medal (IMU).

\textbf{1973:} - Adams Prize (University of Cambridge).

\textbf{1975:} - Elected Fellow of the Royal Society.

\textbf{1980:} - Berwick Prize (London Mathematical Society).

\textbf{2012:} - Honorary Doctorate, University of Nottingham.

\textbf{Elected to:} - Royal Society of London, Fellow. - Hungarian
Academy of Sciences. - American Academy of Arts and Sciences. - Royal
Society of Edinburgh.

\begin{center}\rule{0.5\linewidth}{0.5pt}\end{center}

\section{13. Selected Major
Publications}\label{ch1970alanbaker:selected-major-publications}

\begin{enumerate}
\def\labelenumi{\arabic{enumi}.}
\item
 \textbf{``Linear forms in the logarithms of algebraic numbers I, II,
 III.''} \emph{Mathematika} 13, 204--216, 1966; \emph{Mathematika} 14,
 102--107, 1967; \emph{Mathematika} 14, 220--228, 1967. --- \emph{The
 foundational papers on linear forms in logarithms.}
\item
 \textbf{``The Diophantine equation \(y^2 = ax^3 + bx^2 + cx + d\).''}
 \emph{Journal of the London Mathematical Society} 43, 1--9, 1968. ---
 \emph{Effective solution of Mordell's equation.}
\item
 \textbf{``Contributions to the theory of Diophantine equations I,
 II.''} \emph{Philosophical Transactions of the Royal Society} A 263,
 173--208, 1968. --- \emph{The effective theory of Thue equations.}
\item
 \textbf{``Bounds for solutions of hyperelliptic equations.''}
 \emph{Proceedings of the Cambridge Philosophical Society} 65,
 439--444, 1969. --- \emph{Effective bounds for hyperelliptic
 equations.}
\item
 \textbf{``Linear forms in the logarithms of algebraic numbers IV.''}
 \emph{Mathematika} 15, 204--216, 1968. --- \emph{The \(p\)-adic
 theory.}
\item
 \textbf{``A sharpening of the bounds for linear forms in logarithms I,
 II, III.''} \emph{Acta Arithmetica} 21, 117--129, 1972; \emph{Acta
 Arithmetica} 24, 33--36, 1973; \emph{Acta Arithmetica} 27, 247--252,
 1975. --- \emph{Improvements to the main theorem.}
\item
 \textbf{``Transcendental Number Theory.''} Cambridge University Press,
 1975. --- \emph{The definitive monograph on the subject.}
\item
 \textbf{``Elliptic curves and linear forms in logarithms.''}
 \emph{Journal of Number Theory} 3, 192--210, 1971. ---
 \emph{Application to elliptic curves.}
\item
 \textbf{``On the class number of imaginary quadratic fields.''}
 \emph{Bulletin of the AMS} 72, 912--914, 1966. --- \emph{Determination
 of class number 1 fields.}
\item
 \textbf{``Logarithmic forms and the abc conjecture.''} (with G.
 Wüstholz) In \emph{Number Theory}, 2003. --- \emph{Connections to the
 abc conjecture.}
\end{enumerate}

\begin{center}\rule{0.5\linewidth}{0.5pt}\end{center}

\section{14. References and Further
Reading}\label{ch1970alanbaker:references-and-further-reading}

\subsection{Official Sources}\label{ch1970alanbaker:official-sources}

\begin{itemize}
\tightlist
\item
 IMU Fields Medal 1970:
 https://www.mathunion.org/imu-awards/fields-medal/fields-medals-1970
\item
 Royal Society biographical memoir.
\end{itemize}

\subsection{Key Technical References}\label{ch1970alanbaker:key-technical-references}

\begin{itemize}
\tightlist
\item
 Baker, A. \emph{Transcendental Number Theory}. Cambridge University
 Press, 1975.
\item
 Baker, A. and Wüstholz, G. \emph{Logarithmic Forms and Diophantine
 Geometry}. Cambridge, 2007.
\end{itemize}

\subsection{Textbooks and Surveys}\label{ch1970alanbaker:textbooks-and-surveys}

\begin{itemize}
\tightlist
\item
 Waldschmidt, M. \emph{Diophantine Approximation on Linear Algebraic
 Groups}. Springer, 2000.
\item
 Lang, S. \emph{Elliptic Curves: Diophantine Analysis}. Springer, 1978.
\item
 Shorey, T. N. and Tijdeman, R. \emph{Exponential Diophantine
 Equations}. Cambridge, 1986.
\end{itemize}

\subsection{Exposition and History}\label{ch1970alanbaker:exposition-and-history}

\begin{itemize}
\tightlist
\item
 Masser, D. ``The Work of Alan Baker.'' ICM 1970 Laudatio.
\item
 Schmidt, W. M. \emph{Diophantine Approximation}. Springer, 1980.
\item
 Bertrand, D. ``Baker's Theorem and Transcendence.''
 \emph{L'Enseignement Mathématique}, 1975.
\end{itemize}

\begin{center}\rule{0.5\linewidth}{0.5pt}\end{center}

\emph{(End of report)}
